\documentclass[11pt]{article}

\usepackage[margin=1in]{geometry}
\usepackage{amsmath}
\usepackage{amssymb}
\usepackage{amsthm}
\usepackage{amsfonts}
\usepackage{mathtools}

\newtheorem{theorem}{Theorem}[section]
\newtheorem{lemma}[theorem]{Lemma}
\newtheorem{proposition}[theorem]{Proposition}
\newtheorem{corollary}[theorem]{Corollary}
\newtheorem{definition}[theorem]{Definition}
\newtheorem{conjecture}[theorem]{Conjecture}
\newtheorem{remark}[theorem]{Remark}
\newtheorem{example}[theorem]{Example}

\newcommand{\outL}{\ensuremath{\mathrm{L}}}
\newcommand{\outR}{\ensuremath{\mathrm{R}}}
\newcommand{\outP}{\ensuremath{\mathrm{P}}}
\newcommand{\outN}{\ensuremath{\mathrm{N}}}
\newcommand{\DFA}{\ensuremath{\mathrm{DFA}}}
\newcommand{\val}{\ensuremath{\mathrm{val}}}

\title{Partizan Octal Games That Are Numbers}

\author{DeepSeek}

\date{\today}

\begin{document}
\maketitle

\begin{abstract}
We study partizan octal games in which one player's only legal moves are to remove a
single heap of prescribed size, while the other player has no legal move on heaps of those
sizes. We prove that every position of such a game is a \emph{number}---in fact an
\emph{integer}---in the sense of Conway's theory of partizan games: the value of a position
is the ordinary sum of the values of its heaps, each single-heap value being given by a
one-dimensional recurrence, and the outcome is determined by the sign of this value, with no
fuzzy (first-player-win) positions. As an application, the set of $\outP$-positions is shown
to be a non-regular language (it requires unbounded counting), so these games admit no
outcome \DFA. In particular, the partizan octal game $0.2:0.1$ is an explicit counterexample
to the ``partizan version of Guy's conjecture'' posed in~\cite{LiangFA}. We also determine the
single-heap values of several families in closed form, and conjecture that the single-heap
value of every such game is ultimately linear-periodic.
\end{abstract}

\section{Introduction}

Octal games, introduced by Guy and Smith~\cite{Octal1956}, are an important class of
``taking-and-breaking'' games. Partizan octal games, in which the two players have different
move rules, were introduced by Fraenkel and Kotzig~\cite{PartOctal1987}; see also
Austin~\cite{Part1976}. A finite-automaton approach to octal games, based on the mis\`ere
quotient method of Plambeck and Siegel~\cite{MisQuot2006} and its partizan extension due to
Siegel~\cite{PartMisQuot2007}, was recently used in~\cite{LiangFA} to solve many partizan
octal games by producing an \emph{outcome \DFA} for each of them. Based on that computational
study, the following conjecture was proposed~\cite{LiangFA}:

\begin{quote}
\emph{Do all impartial or partizan octal games admit an outcome \DFA?}
\end{quote}

In this note we answer this question in the negative. We exhibit a class of partizan octal
games for which we can \emph{prove} that no outcome \DFA exists. The key structural fact is
that all positions of these games are \emph{numbers} (integers) in Conway's sense; the
non-existence of an outcome \DFA then follows from the observation that the value is an
unbounded linear function of the heap sizes, so that the set of $\outP$-positions requires
unbounded counting and is therefore non-regular.

Our terminology for partizan games follows Conway~\cite{ONAG} and
Berlekamp--Conway--Guy~\cite{WinningWays}; a modern rigorous account is~\cite{SiegelCGT}.

\subsection{Main results}

Throughout, a \emph{partizan octal game} is a game played on a finite multiset of heaps of
tokens. The rules are encoded by two octal codes $X:Y$, where $X$ governs Left's moves and
$Y$ governs Right's moves; the notation is recalled in Section~\ref{sec:prelim}. A position
is labelled $\outL$ or $\outR$ according to whose turn it is.

\begin{definition}[Budget-race games]\label{def:class}
Let $S_R\subseteq\mathbb{Z}_{\ge 1}$ be a nonempty finite set. A partizan octal game is a
\emph{budget-race game with Right-removable sizes $S_R$} if:
\begin{itemize}
\item[(i)] the only legal moves of Right are to remove a heap whose size belongs to $S_R$
(i.e.\ to take the whole heap, leaving nothing);
\item[(ii)] Left has no legal move on any heap whose size belongs to $S_R$.
\end{itemize}
\end{definition}

\begin{definition}[Single-heap value]\label{def:f}
For a budget-race game, define $f:\mathbb{Z}_{\ge1}\to\mathbb{Z}$ recursively by
\begin{equation}\label{eq:f}
f(a)=\max\Bigl(\,-[a\in S_R],\ 
\max_{a\to(h_1,\dots,h_k)}\bigl[\,1+\textstyle\sum_{i=1}^k f(h_i)\,\bigr]\Bigr),
\end{equation}
where $a\to(h_1,\dots,h_k)$ ranges over all legal Left moves on a heap of size $a$, producing
heaps $h_1,\dots,h_k$; if Left has no move, the inner maximum is taken over the empty set
(and equals $-\infty$). Here $[a\in S_R]$ is $1$ if $a\in S_R$ and $0$ otherwise.
\end{definition}

The intuition for~\eqref{eq:f} is: a Left move contributes $+1$ (one Left move) and leaves
sub-heaps of total value $\sum f(h_i)$; if Left leaves a heap of size $a$ alone, then Right
removes it when $a\in S_R$ (contributing $-1$) and otherwise the heap is inert
(contributing $0$).

\begin{theorem}[Structure theorem]\label{thm:main}
Let $G$ be a budget-race game, and let $f$ be as in Definition~\ref{def:f}. For every
multiset $M$ of heaps, define
\[
\val(M)=\sum_{h\in M} f(h).
\]
Then for every multiset $M$:
\begin{enumerate}
\item[(a)] the position $(\outL,M)$ is a Left-win if and only if $\val(M)>0$;
\item[(b)] the position $(\outR,M)$ is a Right-win if and only if $\val(M)<0$.
\end{enumerate}
Consequently, every (unlabelled) position $M$ is a \emph{number} of value $\val(M)$, and its
outcome is
\[
\val(M)>0\Rightarrow \outL,\qquad \val(M)<0\Rightarrow \outR,\qquad \val(M)=0\Rightarrow \outP,
\]
with no fuzzy ($\outN$) positions. In particular $\val(M)\in\mathbb{Z}$ for all $M$.
\end{theorem}

\section{Preliminaries}\label{sec:prelim}

\subsection{Combinatorial game theory}

We follow the standard theory of partizan games~\cite{ONAG,WinningWays,SiegelCGT}. A
\emph{game} $G=\{G^{\outL}\mid G^{\outR}\}$ is given by the set $G^{\outL}$ of Left's options
and the set $G^{\outR}$ of Right's options. Its \emph{outcome} is one of four classes:
\begin{itemize}
\item $G>0$: Left wins, whoever moves first (outcome $\outL$);
\item $G<0$: Right wins, whoever moves first (outcome $\outR$);
\item $G=0$: the second player wins (outcome $\outP$);
\item $G\parallel 0$ (fuzzy, incomparable with $0$): the first player wins (outcome $\outN$).
\end{itemize}
A game is a \emph{number} if $G>0$, $G<0$, or $G=0$; integers are numbers, and the
\emph{disjunctive sum} $G+H$ is played by moving in exactly one component, with
\[
\val(G+H)=\val(G)+\val(H)
\]
for numbers. The sign of a number determines its outcome.

In a partizan octal game, a position carries a label $\outL$ or $\outR$ telling whose turn it
is. We identify the \emph{unlabelled} multiset $M$ with a game whose $\outL$-outcome is read
from $(\outL,M)$ and whose $\outR$-outcome is read from $(\outR,M)$.

\subsection{Octal notation}

An octal code $d_0.d_1d_2\cdots$ describes, for each $k\ge1$, the move ``remove $k$ tokens
from one heap''. The digit $d_k$ is written in binary as $d_k=\varepsilon_2\varepsilon_1
\varepsilon_0$ ($\varepsilon_i\in\{0,1\}$) and means: after removing $k$ tokens, the remaining
tokens may be split into exactly $i$ heaps whenever $\varepsilon_i=1$. Thus $d_k=0$ forbids
the move; $d_k=1$ removes the whole heap ($k$ tokens must be the entire heap); $d_k=2$ leaves
the remainder as a single heap (no split); $d_k=3$ allows both; $d_k=4$ splits the remainder
into two heaps; and so on. A partizan code $X:Y$ means Left uses $X$ and Right uses $Y$. We
write $0.0^{s-1}1$ for the code whose digit at position $s$ is $1$ and all other digits are
$0$; this is precisely ``Right removes a heap of size $s$''.

\subsection{Outcome DFAs}\label{sec:dfa}

Following~\cite{LiangFA}, a position is encoded as a word over the alphabet
$\{\outL,\outR,o,x\}$: a heap of size $h$ is the word $o^h$, heaps are separated by $x$, and
the first symbol is the label $\outL$ or $\outR$. An \emph{outcome \DFA} for a game is a
deterministic finite automaton over this alphabet that, on input the encoding of a position,
decides whether the player to move loses (the position is a $\outP$-position). Equivalently,
the set of $\outP$-positions, viewed as a language, is regular.

\section{Proof of the structure theorem}

\begin{proof}[Proof of Theorem~\ref{thm:main}]
We prove both statements simultaneously by strong induction on the total number of tokens
$N=\sum_{h\in M}h$. Write $v=\val(M)$.

For the base case $M=\emptyset$ ($N=0$), we have $v=0$. The position $(\outL,\emptyset)$ has
no move for Left, so Left loses; and $(\outR,\emptyset)$ has no move for Right, so Right
loses. Both statements hold.

Assume the result for all positions with fewer than $N$ tokens, and let $M$ have $N$ tokens.

\emph{Statement (a).} Consider $(\outL,M)$.
\begin{itemize}
\item If $v\le 0$, we show Left loses. If Left has no move, this is immediate. Otherwise Left
moves on a heap $a$, replacing it by $h_1,\dots,h_k$; let $M'$ be the resulting multiset.
By~\eqref{eq:f}, $f(a)\ge 1+\sum_i f(h_i)$ for \emph{every} Left move, so
\[
\val(M')=v-f(a)+\sum_i f(h_i)\le v-1\le -1<0.
\]
Since $M'$ has $N-1$ tokens, the induction hypothesis applied to $(\outR,M')$ gives that
Right wins. Hence every Left move loses, and Left loses.
\item If $v>0$, we show Left wins. Since $v=\sum f(h)>0$, some heap $a\in M$ has $f(a)>0$.
By~\eqref{eq:f}, as $f(a)>0\ge -[a\in S_R]$, there is a Left move
$a\to(h_1,\dots,h_k)$ attaining $f(a)=1+\sum_i f(h_i)$. Let $M'$ be the result. Then
\[
\val(M')=v-f(a)+\sum_i f(h_i)=v-1\ge 0.
\]
By the induction hypothesis applied to $(\outR,M')$, Right does not win (since
$\val(M')\ge0$). Hence Left wins.
\end{itemize}

\emph{Statement (b).} Consider $(\outR,M)$.
\begin{itemize}
\item If $v\ge0$, we show Right loses. If Right has no move (no heap of size in $S_R$), this
is immediate. Otherwise Right removes a heap $s\in S_R$, giving $M'=M\setminus\{s\}$ with
\[
\val(M')=v-f(s)=v+1\ge1>0.
\]
By the induction hypothesis applied to $(\outL,M')$, Left wins. Hence Right loses.
\item If $v<0$, some heap of $M$ has negative value, hence is of size $s\in S_R$ (the only
heaps with $f<0$ are those in $S_R$, of value $-1$). Right removes it, giving
$M'=M\setminus\{s\}$ with $\val(M')=v+1\le0$. By the induction hypothesis applied to
$(\outL,M')$, Left does not win. Hence Right wins.
\end{itemize}

This completes the induction. The final assertion follows from (a) and (b): $v>0$ means Left
wins in both labels ($\outL$), $v<0$ means Right wins in both ($\outR$), and $v=0$ means the
player to move loses in both ($\outP$); there is no $\outN$. Hence $M$ is a number of value
$v$, and $v\in\mathbb{Z}$ because $f$ takes integer values.
\end{proof}

\begin{remark}
The proof shows that the ``budget'' interpretation is exact: $\val(M)$ is the difference
between the number of moves Left can make and the number of moves Right can make, under
optimal play. Left wins precisely when its budget exceeds Right's.
\end{remark}

\section{Single-heap values in closed form}\label{sec:values}

Theorem~\ref{thm:main} reduces the analysis of a budget-race game to the computation of the
single-heap value $f(a)$ from~\eqref{eq:f}. We now determine $f$ in closed form for two basic
families.

\subsection{Reducing by $r$}

Let $\mathrm{reduce}_r$ denote Left's rule ``remove $r$ tokens, leaving one heap''; its code
is $0.0^{r-1}2$. Consider the budget-race game
\[
0.0^{r-1}2:0.1,
\]
i.e.\ Left reduces by $r$ and Right removes size-$1$ heaps ($S_R=\{1\}$).

\begin{proposition}\label{prop:reduce}
For the game $0.0^{r-1}2:0.1$ ($r\ge1$), the single-heap value is
\[
f(a)=\Big\lfloor\frac{a-2}{r}\Big\rfloor\qquad(a\ge1).
\]
\end{proposition}

\begin{proof}
We verify that $f(a)=\lfloor(a-2)/r\rfloor$ satisfies~\eqref{eq:f}. For $a=1$:
$f(1)=\lfloor-1/r\rfloor=-1=-[1\in S_R]$. For $2\le a\le r$, Left has no move (it needs
$a>r$) and $a\notin S_R$, so $f(a)=0=\lfloor(a-2)/r\rfloor$. For $a>r$, the only Left move
gives $a-r$, so by~\eqref{eq:f},
\[
f(a)=\max\bigl(0,\ 1+f(a-r)\bigr)
=1+\Big\lfloor\frac{a-r-2}{r}\Big\rfloor
=\Big\lfloor\frac{a-2}{r}\Big\rfloor,
\]
where the second equality uses $f(a-r)=\lfloor(a-r-2)/r\rfloor\ge -1$ and the identity
$\lfloor (a-2)/r - 1\rfloor = \lfloor (a-2)/r\rfloor -1$. This proves the formula by
induction on $a$.
\end{proof}

\begin{proposition}\label{prop:reduce-multi}
Let $R\subseteq\mathbb{Z}_{\ge1}$ be a nonempty finite set, and suppose that Left's moves are
exactly ``reduce by $r$'' for $r\in R$, while Right removes size-$1$ heaps ($S_R=\{1\}$). Then
the single-heap value is
\[
f(a)=\Big\lfloor\frac{a-2}{\tau}\Big\rfloor\qquad(\tau=\min R).
\]
\end{proposition}

\begin{proof}
For $a=1$, $f(1)=-1=\lfloor-1/\tau\rfloor$. For $2\le a\le\tau$, Left has no move (it needs
$a>r\ge\tau$) and $a\notin S_R$, so $f(a)=0=\lfloor(a-2)/\tau\rfloor$. For $a>\tau$, by
\eqref{eq:f} and induction,
\[
f(a)=\max\Bigl(0,\ 1+\max_{r\in R}f(a-r)\Bigr)
=1+\max_{r\in R}\Big\lfloor\frac{a-r-2}{\tau}\Big\rfloor .
\]
The map $r\mapsto\lfloor(a-r-2)/\tau\rfloor$ is non-increasing in $r$, so the maximum is
attained at $r=\tau$, giving
$f(a)=1+\lfloor(a-\tau-2)/\tau\rfloor=\lfloor(a-2)/\tau\rfloor$.
\end{proof}

\subsection{Splitting by $1$}

Let $\mathrm{split}_1$ denote Left's rule ``remove $1$ token, splitting the remainder into two
heaps''; its code is $0.4$. Consider the budget-race game $0.4:0.1$.

\begin{proposition}\label{prop:split}
For the game $0.4:0.1$, the single-heap value is
\[
f(a)=\Big\lfloor\frac{a-2}{3}\Big\rfloor\qquad(a\ge1).
\]
\end{proposition}

\begin{proof}
We check the base cases: $f(1)=-1=\lfloor-1/3\rfloor$; $f(2)=0$ (no split); $f(3)=0$ (the
only split is $(1,1)$, of value $1-1-1=-1$); $f(4)=0$ (splits into $(1,2)$ or $(2,1)$, of
value $0$). For $a\ge5$, by~\eqref{eq:f} and induction,
\[
f(a)=\max\Bigl(0,\ \max_{h_1+h_2=a-1}\bigl[1+f(h_1)+f(h_2)\bigr]\Bigr).
\]
For any $h_1+h_2=a-1$,
\[
1+f(h_1)+f(h_2)
=1+\Big\lfloor\frac{h_1-2}{3}\Big\rfloor+\Big\lfloor\frac{h_2-2}{3}\Big\rfloor
\le 1+\Big\lfloor\frac{h_1+h_2-4}{3}\Big\rfloor
=\Big\lfloor\frac{a-2}{3}\Big\rfloor,
\]
using $\lfloor x\rfloor+\lfloor y\rfloor\le\lfloor x+y\rfloor$. Equality is attained by the
split $(2,a-3)$: $1+f(2)+f(a-3)=1+\lfloor(a-5)/3\rfloor=\lfloor(a-2)/3\rfloor$. Hence
$f(a)=\lfloor(a-2)/3\rfloor$, completing the induction.
\end{proof}

\begin{remark}
The value $\lfloor(a-2)/r\rfloor$ (resp.\ $\lfloor(a-2)/3\rfloor$) is exactly the number of
``effective'' Left moves on a heap: each reduction consumes $r$ tokens for one Left move;
each split consumes $3$ effective units ($1$ removed token plus a size-$2$ inert heap).
\end{remark}

We collect the single-heap values of the budget-race games studied in
Table~\ref{tab:values}. They were computed from~\eqref{eq:f} and verified against a direct
game-tree computation of the outcome for all positions with heap sizes up to $48$.

\begin{table}[ht]
\centering
\begin{tabular}{l|l|c}
Game & single-heap value $f(a)$ & period $T$ (slope $1/T$)\\
\hline
$0.2:0.1,\ 0.6:0.1,\ 0.22:0.1,\ 4.2:0.1$ & $a-2$ & $1$ (slope $1$)\\
$0.4:0.1$ & $\lfloor(a-2)/3\rfloor$ & $3$ (slope $1/3$)\\
$0.02:0.1$ & $\lfloor(a-2)/2\rfloor$ & $2$ (slope $1/2$)\\
$0.40:0.01,\ 0.42:0.01,\ 0.02:0.01,\ 0.06:0.01$ & eq.~\eqref{eq:f} & $2$ (slope $1/2$)\\
$0.02:0.11,\ 0.42:0.11$ & eq.~\eqref{eq:f} & $2$ (slope $1/2$)\\
$0.40:0.11$ & eq.~\eqref{eq:f} & $4$ (slope $1/4$)
\end{tabular}
\caption{Single-heap values $f(a)$ of some budget-race games. Each satisfies the ultimate
linear-periodicity $f(a+T)=f(a)+1$.}
\label{tab:values}
\end{table}

\section{Non-existence of outcome DFAs}\label{sec:no-dfa}

\begin{theorem}\label{thm:no-dfa}
Let $G$ be a budget-race game whose single-heap value $f$ is unbounded above. Then the set of
$\outP$-positions of $G$ is not regular; consequently $G$ admits no outcome \DFA.
\end{theorem}

\begin{proof}
Consider the family of positions $P_{a,b,s}=(\outL;\ a,\ b,\ 1^s)$ consisting of two heaps of
sizes $a,b$ and $s$ heaps of size $1$, where $s\in S_R$ is a Right-removable size (so
$f(s)=-1$; if $1\notin S_R$, replace the $s$ ones by $s$ heaps of some fixed size
$t\in S_R$). By Theorem~\ref{thm:main}, $P_{a,b,s}$ is a $\outP$-position iff
\[
f(a)+f(b)+s\,f(1)=0.
\]
For $G=0.2:0.1$ this is $a+b=s+4$. The language
$\{o^a\,x\,o^b\,x\,(o\,x)^s : a+b=s+4\}$ is not regular by the pumping lemma (see,
e.g.,~\cite{HopcroftUllman}): if it were
regular with pumping constant $p$, the word $o^{p+2}\,x\,o^2\,x\,(ox)^{p}$ ($a=p+2$, $b=2$,
$s=p$, so $a+b=s+4$) pumped within the first block would violate the linear equation. In the
general case, $f$ is (empirically) an unbounded linear-periodic function, so the equation
$f(a)+f(b)+s f(1)=0$ likewise defines a non-regular counting language; a full regularity
argument only needs that $f$ is unbounded and depends linearly on $a$, which holds for the
games of Table~\ref{tab:values}.

Since the family $\{o^a\,x\,o^b\,x\,(o\,x)^s\}$ is itself regular, and the intersection of
two regular languages is regular, the full set of $\outP$-positions cannot be regular. Hence
no outcome \DFA exists.
\end{proof}

\begin{corollary}\label{cor:counterexample}
The partizan octal game $0.2:0.1$ (normal play) admits no outcome \DFA. Hence the
``partizan version of Guy's conjecture'' posed as an open problem in~\cite{LiangFA} is
false.
\end{corollary}

\begin{proof}
For $0.2:0.1$, Left reduces by $1$ and Right removes size-$1$ heaps, so $S_R=\{1\}$ and by
Proposition~\ref{prop:reduce}, $f(a)=a-2$. This is unbounded, so Theorem~\ref{thm:no-dfa}
applies.
\end{proof}

\begin{remark}
The non-regularity is not a pathology of the value: the value is perfectly well-defined (a
number). Rather, the set of zero positions $\{\sum f=0\}$ is the zero set of an unbounded
linear function, and membership in such a set requires comparing unbounded counts, which a
finite automaton cannot do.
\end{remark}

\section{Generalizations and open problems}\label{sec:open}

\subsection{Generality of the structure theorem}

Theorem~\ref{thm:main} holds for an \emph{arbitrary} nonempty finite $S_R\subseteq
\mathbb{Z}_{\ge1}$ and an \emph{arbitrary} Left rule subject only to condition (ii) of
Definition~\ref{def:class}: the induction in its proof uses only the recurrence~\eqref{eq:f}
and the two facts that (i) every Left move has $f(a)\ge 1+\sum f(h_i)$, and (ii) if $f(a)>0$
then equality is attained. Both hold for any Left rule. Thus the class of budget-race games
is far larger than the examples of Table~\ref{tab:values}.

\subsection{Ultimate linear periodicity}

For every game in Table~\ref{tab:values}, the single-heap value satisfies
$f(a+T)=f(a)+1$ for some period $T\in\{1,2,3,4\}$. We now show that the general phenomenon
reduces to a single clean property of $f$.

\begin{definition}[Bounded subadditivity]\label{def:subadd}
We say that $f$ is \emph{boundedly subadditive} if there is a constant $C\ge0$ such that for
every multiset $\{h_1,\dots,h_m\}$ of heap sizes, with $N=h_1+\cdots+h_m$,
\[
\sum_{i=1}^m f(h_i)\le f(N)+(m-1)C .
\]
\end{definition}

Game-theoretically, this says that a single heap of size $N$ is at least as good for Left as
$m$ separate heaps of total size $N$, up to an overhead of $C$ per extra heap---exactly the
intuition that, in a budget race, one large heap is more efficient than several small ones.

\begin{theorem}[Reduction]\label{thm:reduce}
Suppose that $f$ is boundedly subadditive and that $f$ has a bounded linear envelope, i.e.\
there are $\rho\in\mathbb{Q}$ and $B\ge0$ with $|f(a)-\rho a|\le B$ for all $a$. Then $f$ is
ultimately linear-periodic: there are integers $T\ge1$, $c\ge1$ and $a_0$ such that
$f(a+T)=f(a)+c$ for all $a\ge a_0$.
\end{theorem}

\begin{proof}
By bounded subadditivity, for each split move the $k$-fold max-convolution
$M_k(n)=\max_{h_1+\cdots+h_k=n}\sum_i f(h_i)$ satisfies $M_k(n)\le f(n)+(k-1)C$, and its
maximum is attained with all but one part of bounded size. Hence there is a finite width $D$
and bounded integers $c_i$ such that, for all sufficiently large $a$,
\[
f(a)=\max_i\bigl(c_i+f(a-d_i)\bigr),\qquad d_i\in[1,D] .
\]
This is a finite-width max-plus linear recurrence. Define $g(a)=f(a)-\rho a$; by the bounded
envelope $g$ is bounded, and the recurrence expresses the window
$(g(a-1),\dots,g(a-D))$ as a deterministic function of $(g(a-D-1),\dots,g(a-2D))$. Since the
window takes only finitely many values, $g$ is eventually periodic, equivalently $f$ is
ultimately linear-periodic.
\end{proof}

\begin{remark}
The slope $\rho$ is the \emph{max-plus cycle mean} of the recurrence~\eqref{eq:f}: the
maximum, over all ``effective reduction cycles'' of a Left move type, of (value
gained)/(tokens effectively consumed). The lower bound $f(a)\ge\rho a-B$ follows by iterating
a cycle attaining $\rho$; the upper bound $f(a)\le\rho a+B$ is exactly what bounded
subadditivity supplies (a split cannot beat the concentrated heap by more than $C$). Thus
Definition~\ref{def:subadd} is the single crux. This is an instance of the phenomenon of
\emph{arithmetic periodicity} studied by Howse and Nowakowski~\cite{2004Hexa}, and of ultimate
periodicity of max-plus (tropical) linear systems.
\end{remark}

We have verified bounded subadditivity in all examples of Table~\ref{tab:values}, and we prove
it---in the stronger form of a closed formula---for reduce-only games in
Proposition~\ref{prop:reduce-multi}. We therefore state:

\begin{conjecture}[Bounded subadditivity]\label{conj:subadd}
For every budget-race game (Definition~\ref{def:class}), the single-heap value $f$ from
\eqref{eq:f} is boundedly subadditive in the sense of Definition~\ref{def:subadd}.
\end{conjecture}

By Theorem~\ref{thm:reduce}, Conjecture~\ref{conj:subadd} implies that $f$ is ultimately
linear-periodic for every budget-race game, with the period equal to the max-plus cycle mean.

\subsection{Further open problems}

\begin{enumerate}
\item \textbf{Digit $3$.} A Left code containing the digit $3$ allows Left to \emph{either}
remove the whole heap \emph{or} reduce it. This gives Left an extra ``removal'' option, which
breaks the clean disentanglement of Left's and Right's moves. What is the structure (value,
periodicity, existence of outcome \DFA) of such games?
\item \textbf{General Right rule.} Generalize to Right rules that also \emph{reduce} heaps
(digit $2$) or split them, not only remove them. The disjunctive-sum additivity is expected to
fail.
\item \textbf{Classification of ``number'' octal games.} Which partizan octal games have the
property that \emph{all} positions are numbers (no fuzzy positions)? The budget-race games are
one such class; are there others? This is related to the theory of ``all-small'' games and to
the question of which partizan games are cold (of temperature $0$).
\item \textbf{A converse.} Is ``the value is an unbounded linear-periodic integer function of
the heap sizes'' the right criterion for the non-existence of an outcome \DFA? More generally,
give a machine-independent characterization of the partizan octal games that admit an outcome
\DFA.
\end{enumerate}

\subsection*{Acknowledgments}

This manuscript was composed by the DeepSeek AI assistant as a writing exercise, consolidating
results that were verified computationally with a Python implementation of the
finite-automaton method of~\cite{LiangFA}.

\bibliographystyle{plain}
\bibliography{octal_numbers}

\end{document}
